\documentclass{book}
\usepackage[backend=biber,natbib=true,style=authoryear]{biblatex}
\addbibresource{/home/nguyen/1_NQBH/math/bib.bib}
\usepackage[utf8]{inputenc}
\usepackage{float}
\usepackage{graphicx}
\usepackage[colorlinks=true,linkcolor=blue,urlcolor=red,citecolor=magenta]{hyperref}
\usepackage{amsmath,amssymb,amsthm}
\allowdisplaybreaks
\numberwithin{equation}{section}
\newtheorem{assumption}{Assumption}[section]
\newtheorem{lemma}{Lemma}[section]
\newtheorem{corollary}{Corollary}[section]
\newtheorem{definition}{Definition}[section]
\newtheorem{proposition}{Proposition}[section]
\newtheorem{theorem}{Theorem}[section]
\newtheorem{notation}{Notation}[section]
\newtheorem{remark}{Remark}[section]
\newtheorem{example}{Example}[section]
\newtheorem{ques}{Question}[section]
\newtheorem{problem}{Problem}[section]
\newtheorem{conjecture}{Conjecture}[section]
\usepackage[left=0.5in,right=0.5in,top=1.5cm,bottom=1.5cm]{geometry}
\usepackage{fancyhdr}
\pagestyle{fancy}
\fancyhf{}
\addtolength{\headheight}{0pt} % obsolete
\lhead{\scriptsize \chaptername~\thechapter}
\rhead{\scriptsize \nouppercase{\leftmark}} %\nouppercase !
\renewcommand{\chaptermark}[1]{\markboth{#1}{}}
\cfoot{\thepage}
\def\labelitemii{$\circ$}

\title{Kate L. Turabian (2018). \textit{A Manual for Writers of Research Papers, Theses, \& Dissertations}. Chicago Style for Students \& Researchers. 9th ed.}
\author{Collector: Nguyen Quan Ba Hong}
\date{\today}

\begin{document}
\maketitle
\setcounter{secnumdepth}{6}
\setcounter{tocdepth}{6}
\tableofcontents

%------------------------------------------------------------------------------%

\chapter*{A Note to Students}

Known by many as simply “``Turabian'' in honor of its original author, \textit{A Manual for Writers of Research Papers, Theses, and Dissertations} is the authoritative student resource on ``Chicago style.''

This book has helped generations of students successfully research, write, and submit papers in virtually all academic disciplines.

Its guidelines for source citations and style have been condensed and adapted for student writers from another, more comprehensive reference work, \textit{The Chicago Manual of Style}.
\begin{itemize}
    \item Part 1 covers every step of the research and writing process.
    
    It provides practical advice to help you formulate the right questions, read critically, and build arguments.
    
    It also shows you how to draft and revise your papers to strengthen both your arguments and your writing.
    \item Part 2 offers a comprehensive guide to the two methods of Chicago-style source citation, beginning with helpful information on general citation practices in Chap. 15.
    
    In the humanities and most social sciences, you will likely use the notes-bibliography style detailed in Chaps. 16 and 17; in the natural and physical sciences and some social sciences, you will more likely  use the author-date style described in Chaps. 18 and 19.
    \item Part 3 covers Chicago's recommended editorial style, which will help you bring consistency to your writing in matters such as punctuation, capitalization, and abbreviations; this section also includes guidance on     incorporating quotations into your writing and on properly presenting tables and figures.
    \item The appendix presents formatting and submission requirements for theses and dissertations that many academic institutions use as a model, but be sure to follow any local guidelines provided by your institution.
\end{itemize}

\chapter*{Preface}

Students writing research papers, theses, and dissertations in today's colleges and universities inhabit a world filled with digital technologies that were unimagined in 1937 - the year dissertation secretary Kate L. Turabian 1st assembled a booklet of guidelines for student writers at the University of Chicago.

The availability of word-processing software and new digital sources has changed the way students conduct research and write up the results.

But these technologies have not altered the basic task of the student writer: doing well-designed research and presenting it clearly and accurately while following accepted academic standards for citation, style, and format.

%
Turabian's 1937 booklet reflected guidelines found in an already classic resource for writers and editors published by the University of Chicago Press that would ultimately be known as \textit{The Chicago Manual of Style} (\textit{CMOS}).

The Press began distributing Turabian's booklet in 1947 and 1st published the work in book form in 1955, under the title \textit{A Manual for Writers of Term Papers, Theses, and Dissertations}.

Over time, Turabian's book has become a standard reference for students of all levels at universities and colleges across the country.

Turabian died in 1987 at age 94, a few months after publication of the book's 5th edition.

%
Beginning with that edition, members of the Press editorial staff have carried out the revisions to the chapters on source citation, style, and paper format.

For the 7th edition (2007), Wayne C. Booth, Gregory G. Colomb, and Joseph M. Williams expanded the focus of the book by adding extensive new material adapted from their book \textit{The Craft of Research}, also published by the University of Chicago Press and now in its 4th edition (2016).

Among the new topics covered in their chapters were the nature of research, finding and engaging sources, taking notes, developing an argument, drafting and revising, and presenting evidence in tables and figures.

Following the deaths of this remarkable trio of authors, whose collective voice will always animate this work, Joseph Bizup and William T. FitzGerald have with this edition assumed the mantle of revising their chapters for a new generation of students, as they had previously for \textit{The Craft of Research}.

%
Part 1, now aligned with the most recent edition of that book, incorporates updated advice for writers and responds to recent developments in information literacy, including the use of digital materials.

Part 2 offers a comprehensive guide to the 2 Chicago styles of source citation - the notes-bibliography format used widely in the humanities and most social sciences and the author-date format favored in many of the sciences and some social sciences.

Thoroughly updated guidance related to online citation practices has been supplemented throughout by new examples featuring the types of sources students are most likely to consult.

Part 3 addresses matters of spelling, punctuation, abbreviation, and treatment of numbers, names, special terms, and titles of works.

The final 2 chapters in this section treat the mechanics of using quotations and graphics (tables and figures), topics that are discussed from a rhetorical perspective in part 1.

Both parts 2 and 3 have been updated for this edition
in accordance with the 17th edition (2017) of \textit{The Chicago Manual of Style}.

The recommendations in this manual in some instances diverge from \textit{CMOS} in small ways, to better suit the requirements of academic papers as opposed to published works.

%
The appendix presents guidelines for paper format and submission that have become the primary authority for dissertation offices throughout the United States.

These guidelines have been updated to reflect the nearly universal electronic submission of papers and to feature new examples from recently published dissertations.

This appendix is intended primarily for students writing PhD dissertations and master's and undergraduate theses, but the sections on format requirements and electronic file preparation will also aid those writing class papers.

An extensive bibliography, organized by subject area and fully updated, lists sources for research and style issues specific to various disciplines.

%
The guidelines in this manual offer practical solutions to a wide range of issues encountered by student writers, but they may be supplemented - or even overruled - by the conventions of specific disciplines or the preferences of particular institutions, departments, or instructors.

All of the chapters on style and format remind students to review the requirements of their university, department, or instructor, which take precedence over the guidelines presented in this book.

%
Updating a book that has been used by millions of students over 80 years is no small task, and many people participated in preparing this 9th edition.

The Press staff welcomed Bizup and FitzGerald to their new role in revising part 1.

Russell David Harper, the principal reviser of the 16th and 17th editions of \textit{CMOS}, revised parts 2 and 3 and the appendix.

Several recent PhD recipients from the University of Chicago allowed the use of excerpts from their dissertations in the appendix, where they are credited individually.

%
Within the Press, the project was developed under the guidance of editors Mary E. Laur and David Morrow, editorial director Christie Henry, and editorial associates Rachel Kelly and Susan Zakin.

Lucy Johnson and Kristin Zodrow offered additional research support.

Ruth Goring edited the manuscript, June Sawyers proofread the pages, and James Curtis prepared the index.

Michael Brehm provided the design, while Joseph Claude supervised the production.

Carol Kasper, Jennifer Ringblom, Lauren Salas, and Carol Fisher Saller brought the final product to market.
\begin{flushright}
    The University of Chicago Press Editorial Staff
\end{flushright}

\part{Research \& Writing}

Wayne C. Booth, Gregory G. Colomb, Joseph M. Williams, Joseph Bizup, and William T. FitzGerald

\chapter*{Overview of Part I}

We know how daunting it can feel to start a substantial research project, whether it's a doctoral dissertation, a master's or senior thesis, or just a long class paper.

But you can handle any project if you \textit{break it into its parts, then work on them 1 step at a time}.

Part I of this book shows you how.

%
We 1st discuss the aims of research and what readers will expect of any research paper (a term we use broadly to refer to all varieties of research-based writing).

We then focus on how to find a research question and problem whose answer is worth your time and your readers' attention; how to find and use information from sources to back up your answer; and then \textit{how to plan, draft, and revise your paper} so your readers will see that \textit{your answer is based on sound reasoning and reliable evidence}.

%
Several themes run through this part.
\begin{itemize}
    \item You can't plunge into a project blindly; you must plan it, then keep the whole process in mind as you take each step.
    
    So think big, but break the process down into small goals that you can meet one at a time.
    \item Your best research will begin with a question that \textit{you} want to answer.
    
    But you must then imagine readers asking questions of their own: \textit{So what if you don't answer it?}
    
    \textit{Why should I care?}
    \item From the outset, you should try to write every day, not just to record the content of your sources but to clarify what you think of them.
    
    You should also write down your own developing ideas to get them out of the cozy warmth of your head and into the cold light of day, where you can see if they still make sense.
    
    You probably won't use much of this writing in your final draft, but it is essential preparation for it.
    \item No matter how carefully you do your research, readers will judge it by how well you present it, so you \textit{must know what they will look for in a clearly written paper that earns their respect}.
\end{itemize}
If you're an advanced researcher, skim Chaps. 1--4.

You will see there much that's familiar; but if you're also teaching, it may help you \textit{explain what you know to your students more effectively}.

Many experienced researchers tell us that Chaps 5--12 have helped them not only to \textit{explain to others how to conduct and report research}, but also \textit{to draft and revise their own writing more quickly and effectively}.

%
If you're just starting your career in research, you'll find every chapter of part 1 useful.

Skim it all for an overview of the process; then as you work through your project, reread chapters relevant to your immediate task.

%
You may feel that the steps described here are too many to remember, but you can manage them if you take them one at a time, and as you do more research they'll become habits of mind.

Don't think, however, that you must follow these steps in exactly the order we present them.

Researchers regularly think ahead to future steps as they work through earlier ones and revisit earlier steps as they deal with a later one.

(That explains why we so often refer you ahead to anticipate a later stage in the process and back to revisit an earlier one.)

And even the most systematic researcher has unexpected insights that send her off in a new direction.

\textit{Work from a plan, but be ready to depart from it, even to discard it for a new one}.

%
If you're a very new researcher, you may also think that some matters we discuss are beyond your immediate needs.

We know that a 10-page class paper differs from a dissertation.

But both require a kind of thinking that even the newest researcher can start practicing.

You begin your journey toward full competence when you not only know what lies ahead but also can start practicing the skills that experienced researchers began to learn when they were where you are now.

%
No book can prepare you for every aspect of every research project.

And this one won't help you with the specific methodologies in fields such as psychology, economics, and philosophy, much less physics, chemistry, and biology.

Nor does it tell you how to adapt what you learn about academic research to business or professional settings.

%
But it does provide \textit{an overview of the processes} and \textit{habits of mind that underlie all researc}h, wherever it's done, and of the plans you must make to \textit{assemble a paper, draft it, and revise it}.

With that knowledge and with help from your teachers, you'll come to \textit{feel in control of your projects}, not intimidated by them, and eventually you'll learn to manage even the most complex projects on your own, in both the academic and the professional worlds.

%
\textit{The 1st step in learning the skills of sound research is to understand how experienced researchers think about its aims}.

%------------------------------------------------------------------------------%

\chapter{What Research Is \& How Researchers Think about It}

Whenever we read about a scientific breakthrough or a crisis in world affairs, we benefit from the research of others, who likewise benefited from the research of countless others before them.

When we walk into a library, we are surrounded by more than 25 centuries of research.

When we go on the internet, we can read the work of millions of researchers who have posed questions beyond number, gathered untold amounts of information from the research of others to answer them, and then shared their answers with the rest of us.

We can carry on their work by asking and, we hope, answering new questions in turn.

Governments spend billions on research, businesses even more.

Research goes on in laboratories and libraries, in jungles and ocean depths, in caves and in outer space, in offices and, in the information age, even in our own homes.

\textit{Research is in fact the world's biggest industry}.

%
So what, exactly, is it?

\section{What Research Is}
You already have a basic understanding of research: \textit{answering a question by obtaining information}.

In this sense, research can be as simple as choosing a new phone or as complex as discovering the origin of life.

In this book we use \textit{research} in a specific way to mean \textit{a process of systematic inquiry to answer a question that not only the researcher but also others want to solve}.

Research thus includes the steps involved in presenting or reporting it.

\textit{To be a true researcher}, as we are using the term, you \textit{must share your findings and conclusions with others}.

%
If you are new to research, you may think that your paper will add little to the world's knowledge.

But done well, it will add a lot to \textit{your} knowledge and to your ability to communicate that knowledge.

As you \textit{learn to do your own research}, you \textit{also learn to use and judge that of others}.

In every profession, researchers must read and evaluate the work of others before they make a decision.

This is a job you will do better after you have learned how others judge yours.

%
This book focuses on research in the academic world, but every day we read or hear about research that affects our lives.

Often we get news of research secondhand, and it can be difficult to know what reasoning and evidence support a claim.

But research doesn't ask for our blind trust or that we accept something on the basis of authority.

It invites readers to think critically about evidence and reasoning.

%
That is how \textit{research-based writing} differs from other kinds of persuasive writing: it must rest on shared facts that readers accept as truths independent of your feelings and beliefs.

Your readers must be able to follow your reasoning from evidence they accept to the claim you draw from it.

\textit{Your success as a researcher thus depends} not just on \textit{how well you gather and analyze data} but also on \textit{how clearly you report your reasoning} so that your readers can test and judge it before making your claims part of their knowledge and understanding.

\section{How Researchers Think about Their Aims}
All researchers collect information, what we're calling \textit{data}.

But researchers do not merely gather facts on a topic - \textit{stories about the Battle of the Alamo}, for example.

They look for specific data to test and support an answer to a question that their topic inspired them to ask, such as \textit{Why has the Alamo story become a national legend?}

In doing so, they also imagine a community of readers who they believe will share their interest and help them test and support an answer to that question.

%
Experienced researchers, however, know that they must do more than convince us that their answer is sound.

They must also show us \textit{why their question was worth asking, how its answer helps us understand some bigger issue in a new way}.

\textit{If we can figure out why the Alamo story has become a national legend, we might then answer a larger question: how have regional myths shaped the American character?}

%
You can \textit{judge how closely your thinking tracks that of an experienced researcher by describing your project} in a sentence like this:
\begin{enumerate}
    \item Topic: I am working on X (\textit{stories about the Battle of the Alamo})
    \begin{itemize}
        \item[2.] Question: because I want to find out Y (\textit{why its story became a national legend})
        \begin{itemize}
            \item[3.] Significance: so that I can help others understand Z (\textit{how such regional myths have shaped the American character}).
        \end{itemize}
    \end{itemize}
\end{enumerate}
That sentence is worth a close look, because it describes not just the progress of your research but your personal growth as a researcher.
\begin{enumerate}
    \item Topic: ``I am working on X$\ldots$'': Those new to research often begin with a simple topic like \textit{the Battle of the Alamo}.
    
    But too often they stop there, with nothing but a broad topic to guide their work.
    
    Beginning this way, they may pile up dozens or hundreds of notes but then can't decide what data to keep or discard.
    
    When it comes time to write, their papers become ``data dumps'' that leave readers wondering what all those data add up to.
    \item Question: ``$\ldots$ because I want to find out Y $\ldots$'': More experienced researchers begin not just with a topic but with a \textit{research question}, e.g. \textit{Why has the story of the Alamo become a national legend?}
    
    They know that readers will think their data add up to something only when they serve as evidence to support an answer.
    
    Indeed, \textit{only with a question can a researcher know what information to look for and, once obtained, what to keep} - and not just data that support a particular answer but also data that test or discredit it.
    
    With sufficient evidence to support an answer, a researcher can respond to data that seem to contradict it.
    
    In writing a paper, the researcher tests that answer and invites others to test it too.
    \item Significance: ``$\ldots$ so that I can help others understand Z'': The best researchers understand that readers want to know not only that \textit{an answer is sound} but also \textit{why the question is worth asking: So what?}
    
    \textit{Why should I care why the Alamo story has become a national legend?}
    
    Think of it this way: what will be lost if you \textit{don't} answer your question?
    
    Your answer might be \textit{Nothing}.
    
    \textit{I just want to know}.
    
    Good enough to start but not to finish, because eventually your readers will want an answer beyond \textit{Just curious}.
\end{enumerate}
Answering \textit{So what?} is tough for all researchers, beginning and experienced alike, because when you only have a question stemming from a topic of personal interest, it's hard to predict whether others will find its answer significant.

Some researchers therefore work backwards: they begin not by following their own curiosity but by crafting questions with implications for bigger ones that others in their field already care about.

But many researchers, including us, find that they \textit{cannot address that 3rd step until they finish a 1st draft}.

So it's fine to \textit{begin} your research without being able to answer \textit{So what?}, and if you are a student, your teacher may even let you skip that last step.

But if you are doing advanced research, you \textit{must} take it, because your answer to \textit{So what?} is what makes your research
matter to others.

%
In short, not all questions are equally good.

We might ask how many cats slept in the Alamo the night before the battle, but so what if we find out?

It is hard to see how an answer would help us think about any larger issue worth understanding, so it's a question that's probably not worth asking (though as we'll see, we could be wrong about that).

%
How good a question is depends on its significance to some community of readers.

Exactly what community depends on your field but also on how you frame your research.

You can try to expand your potential readership by connecting Z to even broader questions: \textit{And if we can understand what has shaped the American character, we might understand better who Americans think they are}.

\textit{And when we know that, we might better understand why others in the world judge them as they do}.

Now perhaps political scientists will be as interested in this research as historians.

On the other hand, if you try to widen your audience too much, you risk losing it altogether.

Sometimes it's better to address a smaller community of specialists.

%
We can't tell you the right choice, but we can tell you 2 wrong ones: trying to interest everyone (some people just won't care no matter how you frame your research) or not trying to interest anyone at all.

\section{Conversing with Your Readers}
When you can explain the \textit{significance} of your research, you \textit{enter into a kind of conversation with your research community}.

Some people, when they think of research, imagine a lone scholar or scientist in a hushed library or lab.

\textit{But no places are more crowded with the presence of others than these}.

When you read a book or an article or a report, you  silently converse with its authors - and through them with everyone else they have read.

In fact, every time you go to a written source for information, you join a conversation between writers and readers that began millennia ago.

And when you report your own research, you add your voice and hope that other voices will respond to you, so that you can in turn respond to them.

And so it goes.

%
Experienced researchers understand that they are participating in such conversations and that genuine research must matter not only to the researcher but also to others.

That is why our formula - \textit{I am working on X to find out Y so that others can better understand Z} - is so powerful: \textit{because it makes informing others the end of research}.

%
But these silent conversations differ from the face-to-face conversations we have every day.

We can judge how well everyday conversations are going as we have them, and we can adjust our statements and behavior to repair mistakes and misunderstandings as they occur.

But in writing we don't have that opportunity: readers have to \textit{imagine} writers in conversation with one another, as well as with themselves, and writers have to imagine their readers and their relationship to them.

In other words, writers have to offer readers a social contract: \textit{I'll play my part if you play yours}.

%
Doing this is 1 of the toughest tasks for beginning researchers: get that relationship wrong and your readers will think you are naive or, worse, won't read your work at all.

Too many beginning researchers offer their readers a relationship that caricatures a bad classroom: \textit{Teacher, I know less than you}.

\textit{So my role is to show you how many facts I can dig up}.

\textit{Yours is to say whether I've found enough to give me a good grade}.

Do that and you turn your project into a pointless drill, casting yourself in a role exactly opposite to that of a true researcher.

In true research, you must switch the roles of student and teacher.

You must imagine a relationship that goes beyond \textit{Here are some facts I've dug up about 14th-century Tibetan weaving}.

\textit{Are they enough of the right ones?}

%
There are 3 better reasons to share what you've found.

You could say to your reader, \textit{Here is some information that you may find interesting}.

This offer assumes, of course, that your reader wants to know.

You could also say not just \textit{Here is something that should interest you} but \textit{Here is something that will help you remedy a situation that troubles you}.

People do this kind of research every day in business, government, and the professions when they try to figure out how to address problems ranging from insomnia to falling profits to climate change.

In Chap. 2 we call such situations and their consequences \textit{practical problems}.

When academic researchers address such practical problems, we say they are doing \textit{applied} research.

Most commonly, though, academic researchers do \textit{pure} research that addresses what we call \textit{conceptual problems} - i.e., \textit{not troubling situations in the world but the limitations of our understanding of it} (again see Chap. 2).

In this case, you say to your readers, \textit{Here is something that will help you better understand something you care about}.

When you make this last sort of appeal, you imagine your readers as a community of receptive but also skeptical colleagues who are open to learning from you and even changing their minds - if you can make the case.

%
We now understand the \textit{goal of research}, at least in its pure form: \textit{it is not to have the last word but to keep the conversation going}.

\textit{The best questions are those whose answers raise several more}.

\textit{When that happens, everyone in the research community benefits}.

%------------------------------------------------------------------------------%

\chapter{Defining a Project: \textit{Topic, Question, Problem, Working Hypothesis}}

A research project begins well before you search the internet or head for the library and continues long after you have collected all the data you think you need.

\textit{Every project involves countless specific tasks, so it is easy to get overwhelmed}.

But in all research projects, you have just 5 general aims:
\begin{enumerate}
    \item Ask a question worth answering.
    \item Find an answer that you can support with good reasons.
    \item Find good data that you can use as reliable evidence to support your reasons.
    \item Draft an argument that makes a good case for your answer.
    \item Revise that draft until readers will think you met the 1st 4 goals.
\end{enumerate}
You might even post those 5 goals in your workspace.

%
Research projects would be much easier if we could march straight through these steps.

But you will discover (if you have not already) that the \fbox{\textit{research process is not so straightforward}.}

Each task overlaps with others, and frequently you must go back to an earlier one.

The truth is, \fbox{\textit{research is messy and unpredictable}.}

But that's also what makes it \textit{exciting} and \textit{ultimately rewarding}.

\section{Find a Question in Your Topic}
Researchers begin projects in different ways.

Many experienced researchers begin with a question that others in their field want to answer: \textit{What caused the extinction of most large North American mammals?}

Others begin with \textit{just basic curiosity, a vague intellectual itch that they have to scratch}.

They might not know what puzzles them about a topic, but they're willing to spend time to find out whether that topic can yield a question worth answering.

%
They realize, moreover, that \textit{the best research question is not one whose answer they want to know just for its own sake; it is one that helps them and others understand some larger issue}.

E.g., if we knew why North American sloths disappeared, we might be able to answer a larger question that puzzles many historical anthropologists: \textit{Did early Native Americans live in harmony with nature, as some believe, or did they hunt its largest creatures to extinction?}

\textit{And if we knew that, then we might also understand}$\ldots$ (\textit{So what?} again. See 1.2.)

%
Then there are those questions that just pop into a researcher's mind with no hint of where they'll lead, sometimes about matters so seemingly trivial that only the researcher thinks they're worth answering: \textit{Why does a coffee spill dry up in the form of a ring?}

Such a question might lead nowhere, but you can't know that until you see its answer.

In fact, the scientist puzzled by coffee rings made discoveries about the behavior of fluids that others in his field thought important - and that paint manufacturers found valuable.

\textit{If you cultivate the ability to see what's odd in the commonplace, you'll never lack for research projects as either a student or a professional}.

%
If you already have a focused topic, you might skip to 2.1.3 and begin asking questions about it.

If you already have some questions, skip to 2.1.4 to test them using the criteria listed there.

Otherwise, here's a plan to help you search for a topic.

\subsection{Search Your Interests}
Beginning researchers often find it hard to pick a topic or believe they lack the expertise to research a topic they have.

But a research topic is an interest stated specifically enough for you to imagine \textit{becoming} a local expert on it.

That doesn't mean you \textit{already} know a lot about it or that you'll know more about it than others, including a teacher or advisor.

\fbox{\textit{You just want to know more about it than you do now}.}

%
If you can work on any topic, we offer only a cliché: \textit{start with what interests you}.

Ask these questions:
\begin{enumerate}
    \item What special interests do you have - chess, old comic books, scouting?
    
    The less common, the better.
    
    Choose one and investigate something about it that you don't know.
    \item Where would you like to travel?
    
    Find out all you can about your destination.
    
    What particular aspect surprises you or makes you want to learn more?
    \item Can you find an online discussion list or social media page focused on issues that interest you?
    \item Visit a museum or a ``virtual museum'' on the internet with exhibitions that appeal to you.
    
    What catches your interest that you would like to know more about?
    \item Have you taken positions on issues in your field or in debates with others but found that you couldn't back up your views with good reasons and evidence?
    \item What issues in your field do people outside your field misunderstand?
    \item What topic is your instructor or advisor working on?
    
    Would she like you to explore a part of it?
    
    Don't be too shy to ask.
    \item Does your library have rich resources in some field?
    
    Ask your instructor or a librarian.
    \item What intrigues you in your reading?
    
    What connections do you see among different things you are reading?
    \item What other courses will you take in your field or out of it?
    
    Find a textbook and skim it for study questions.
    \item If you have a job in mind, what kind of writing might help you get it?
    
    Employers often ask for samples of an applicant's work.
\end{enumerate}
Once you have a list of possible topics, choose 1 or 2 that interest you most and explore their research potential.

Sometimes beginning researchers choose a topic because they already know what they want to say about it, even before they've done any research.

That's a mistake: \textit{the best topics provoke good questions; the worst come with ready-made answers}.

To gauge a topic's potential, do these things:
\begin{itemize}
    \item In the library, look up your topic in a general guide such as \textit{CQ Researcher} and skim the subheadings.
    
    In an online database such as Academic Search Premier, you can explore your topic through subject terms.
    
    If you have a narrower focus, you can do the same with specialized guides such as \textit{Women's Studies International}.
    
    At most libraries today, such guides are found online.
    \item On the internet, google your topic, but don't surf indiscriminately.
    
    Look 1st for websites that are roughly like the sources you would find in a library, such as online encyclopedias.
    
    Read the entries on your general topic, and then copy their lists of references for a closer look.
    
    Few experienced researchers trust Wikipedia as a reliable source to cite as evidence, but most would use the site to find ideas and more specific sources.
    \item Finally, think ahead: \textit{you may be in for a long relationship with your topic, so be sure it interests you enough to get you through the inevitable rocky stretches}.
\end{itemize}

\subsection{Make Your Topic Manageable}
If you pick a topic that sounds like an encyclopedia entry - \textit{bridges, birds, masks} - you'll find so many sources that you could spend a lifetime reading them.

You must carve out of your topic a manageable piece.

Before you start searching, limit your topic to reflect a special interest in it: What is it about, say, masks that made you choose them?

What particular aspect of them interests or puzzles you?

Think about your topic in a context that you know something about, and then add words and phrases to reflect that knowledge:
\begin{example}
    masks in religious ceremonies
    
    masks as symbols in Hopi religious ceremonies
    
    mudhead masks as symbols of sky spirits in Hopi fertility ceremonies
\end{example}
You might not be able to focus your topic until after you start reading about it.

\textit{That takes time, so start early} (you can do much of this preliminary work online):
\begin{itemize}
    \item Begin with an overview of your topic in a general encyclopedia (in the bibliography, see items in category 2 in the general sources); then read about it in a specialized one (see items in category 2 in your field).
    \item Skim a survey of your topic (encyclopedia entries usually cite a few).
    \item Skim subheads under your topic in an annual bibliography in your field (in the bibliography, see items in category 4 in your field).
    
    That will also give you a start on a reading list.
    \item Search the internet for the topic (but evaluate the reliability of what you find; see 3.3.2).
\end{itemize}
Especially useful are topics that spark debate: \textit{Fisher claims that Halloween masks reveal children's archetypal fears, but do they?}

\textit{Even if you can't resolve the debate, you can learn how such debates are conducted} (for more on this, see 3.1.2).

\subsection{Question Your Topic}
Once they have a focused topic, many new researchers start plowing through all the sources they can find, taking notes on everything they read.

They then dump it all into a report with little sense of purpose or direction.

Experienced researchers, however, document information not for its own sake but to support an answer to a question they (and they hope their readers) think worth asking.

So the best way to begin working on a focused topic is to pose questions that direct you to just the information you need to answer them.

%
Do this not just once, early on, but throughout your project.

Ask questions as you read, especially \textit{how} and \textit{why} (see also 4.1.1--4.1.2).

Try the following kinds of questions (the categories are loose and overlap, so don't worry about keeping them distinct).
\begin{enumerate}
    \item Ask how the topic fits into a larger context (historical, social, cultural,
    geographic, functional, economic, and so on):
    \begin{itemize}
        \item How does your topic fit into a larger story?
        
        \textit{What came before masks?}
        
        \textit{How did masks come into being?}
        
        \textit{Why?}
        
        \textit{What changes have they caused in other parts of their social or geographic setting?}
        
        \textit{How and why did that happen?}
        
        \textit{Why have masks become a part of Halloween?}
        
        \textit{How and why have masks helped make Halloween the biggest American holiday after Christmas?}
        \item How is your topic a functioning part of a larger system?
        
        \textit{How do masks reflect the values of specific societies and cultures?}
        
        \textit{What roles do masks play in Hopi dances?}
        
        \textit{In scary movies?}
        
        \textit{In masquerade parties?}
        
        \textit{For what purposes are masks used other than disguise?}
        
        \textit{How has the booming market for kachina masks influenced traditional designs?}
        \item How does your topic compare to and contrast with other topics like it?
        
        \textit{How do masks in Native American ceremonies differ from those in Africa?}
        
        \textit{What do Halloween masks have to do with Mardi Gras masks?}
        
        \textit{How are masks and cosmetic surgery alike?}
    \end{itemize}
    \item Ask questions about the nature of the thing itself, as an independent entity:
    \begin{itemize}
        \item How has your topic changed through time?
        
        Why?
        
        What is its future?
        
        \textit{How have Halloween masks changed?}
        
        \textit{Why?}
        
        \textit{How have Native American masks changed?}
        
        \textit{Why?}
        \item How do the parts of your topic fit together as a system?
        
        \textit{What parts of a mask are most significant in Hopi ceremonies?}
        
        \textit{Why?}
        
        \textit{Why do some masks cover only the eyes?}
        
        \textit{Why do so few masks cover just the bottom half of the face?}
        \item How many different categories of your topic are there?
        
        \textit{What are the different kinds of Halloween masks?}
        
        \textit{What are the different qualities of masks?}
        
        \textit{What are the different functions of Halloween masks?}
    \end{itemize}
    \item Turn positive questions into a negative ones: \textit{Why have masks \emph{not} become a part of Christmas?}
    
    \textit{How do Native American masks \emph{not} differ from those in Africa?}
    
    \textit{What parts of masks are typically \emph{not} significant in religious ceremonies?}
    \item Ask speculative questions: \textit{Why are masks common in African religions but not in Western ones?}
    
    \textit{Why are children more comfortable wearing Halloween masks than are most adults?}
    
    \textit{Why don't hunters in camouflage wear masks?}
    \item Ask \textit{What if?} questions: how would things be different if your topic never existed, disappeared, or were put into a new context?
    
    \textit{What if no one ever  wore masks except for safety reasons?}
    
    \textit{What if everyone wore masks in public?}
    
    \textit{What if movies and TV were like Greek plays and all the actors wore masks?}
    
    \textit{What if it were customary to wear masks on blind dates?}
    
    \textit{In marriage ceremonies?}
    
    \textit{At funerals?}
    \item Ask questions that reflect disagreements with a source: if a source makes a claim you think is only weakly supported or even wrong, make that disagreement a question (see also 4.1.2).
    
    \textit{Martinez claims that carnival masks uniquely allow wearers to escape social norms}.
    
    \textit{But I think religious masks also allow wearers to escape from the material realm to the spiritual}.
    
    \textit{Is there a larger pattern of all masks creating a sense of alternative forms of social or spiritual life?}
    \item Ask questions that build on agreement: if a source offers a claim you
    think is persuasive, ask questions that extend its reach (see also 4.1.1).
    
    \textit{Elias shows that masked balls became popular in 18th-century London in response to anxiety about social mobility}.
    
    \textit{Is the same anxiety responsible for similar developments in other European capitals?}
    
    You can also ask a question that supports the same claim with additional evidence.
    
    \textit{Elias supports his claim about masked balls entirely with published sources}.
    
    \textit{Is it also supported by evidence from unpublished sources such as letters and diaries?}
    \item Ask questions analogous to those that others have asked about similar topics.
    
    \textit{Smith analyzed the Battle of Gettysburg from an economic point of view}.
    
    \textit{What would an economic analysis of the Battle of the Alamo turn up?}
    \item Look for questions that other researchers pose but don't answer.
    
    Many journal articles end with a paragraph or two about open questions, ideas for more research, and so on.
    
    You might not be able to do all the research they suggest, but you might carve out a piece of it.
    \item Find a professional discussion forum on your topic, then ``lurk,'' just reading the exchanges to understand the kinds of questions being asked.
    
    If you can't find one using a search engine, ask a teacher or visit websites of professional organizations in your field.
    
    Look for questions that spark your interest.
    
    If questions from students are welcomed, you can even post one yourself, so long as it is very specific and narrowly focused.
\end{enumerate}

\subsection{Evaluate Your Questions}
After asking all the questions you can think of, evaluate them.

Not all questions are equally good.

Look for questions whose answers might make you (and your readers) think about your topic in a new way.

Avoid questions like these:
\begin{itemize}
    \item Their answers are settled fact that you could just look up.
    
    \textit{What was Audre Lorde's 1st published poem?}
    
    Questions that ask \textit{how} and \textit{why} call for interpretations, not just the discovery of facts.
    
    That's why they invite deeper thinking than questions beginning \textit{who, what, when}, or \textit{where}, and deeper thinking leads to more interesting answers.
    \item Their answers can't be plausibly disproved.
    
    \textit{How important are masks in Inuit culture?}
    
    The answer is obvious: \textit{Very}.
    
    If you can't imagine disproving a claim, then proving it is pointless.
    
    (On the other hand, world-class reputations have been won by those who questioned a claim that seemed self-evidently true - for instance, that the sun circled the earth - and dared to disprove it.)
    \item Their answers would be merely speculative.
    
    \textit{Would church services be as well attended if the congregation all wore masks?}
    
    If you can't imagine finding data that would settle the question, it's not a question you can answer.
    \item Their answers are dead ends.
    
    \textit{How many black cats slept in the Alamo the night before the battle?}
    
    It's hard to see how an answer would help us think about any larger issue worth  understanding better, so the question is probably not worth asking.
    \item Their answers require different capacities from the ones you have.
    
    \textit{How do Japanese translations of \emph{The Great Gatsby} treat early 20th-century America?}
    
    If you can't read Japanese, this question is not for you to answer.
    \item Their answers require more or different resources - materials, technology, money, especially time - than you have.
    
    \textit{How is childhood represented in the Victorian novel?}
    
    Can you read enough of them in the time you have to arrive at a reasonable answer?
\end{itemize}
%
Don't reject a question because you think someone must already have asked it.

Until you know, pursue its answer as if you asked 1st.

Even if someone has answered it, you might come up with a better answer or at least one with a new slant.

In fact, in the humanities and social sciences the best questions usually have more than 1 good answer.

You can also organize your project around comparing and contrasting competing answers and supporting the best one (see 6.2.5).

%
The point is to find a question that \textit{you} want to answer.

\textit{Too many students, both graduate and undergraduate, think that the aim of education is to learn settled answers to someone else's questions}.

It's not.

\textit{It is to find your own answers to your own questions}.

To do that, you must learn to wonder about things, to let them puzzle you - particularly things that seem commonplace.

\section{Understanding Research Problems}
In Chap. 1 we gave you a formula that expresses how experienced researchers think about their work:
\begin{enumerate}
    \item Topic: I am working on X (\textit{stories about the Battle of the Alamo})
    \begin{itemize}
        \item[2.] Question: because I want to find out Y (\textit{why its story became a national legend})
        \begin{itemize}
            \item[3.] Significance: so that I can help others understand Z (\textit{how such regional myths have shaped the American character}).
        \end{itemize}
    \end{itemize}
\end{enumerate}
When you can state that significance from the point of view of your readers, you have more than a question: you have posed a \textit{research problem} that they recognize needs a solution.

%
Among researchers, the term \textit{problem} has a special meaning that sometimes confuses beginners.

In our everyday world, a problem is something we try to avoid.

But in academic research, a problem is something we seek out, even invent.

Indeed, without a problem to work on, a researcher is out of work.

%
Experienced researchers often talk about their problems in shorthand.

When asked what they are working on, they often answer with what sounds like a general topic: \textit{adult measles, mating calls of Wyoming elk}.

As a result, beginners may think that having a topic to read about is the same thing as having a problem to solve.

But \textit{without a specific question to answer and a reason to find that answer significant, researchers have no way of knowing when they have enough}.

So they can be tempted to throw in everything just to be safe.

%
To avoid the judgment that your paper is just a \textit{data dump}, you need a problem, one that focuses on finding just those data that will help you solve it.

To find one, you need to know how problems work.

\subsection{Understanding Practical and Conceptual Problems}
There are 2 kinds of research problems: \textit{practical} and \textit{conceptual}.

Each of them has a 2-part structure:
\begin{itemize}
    \item a situation or \textit{condition}, and
    \item undesirable \textit{costs} or \textit{consequences} caused by that condition.
\end{itemize}
\fbox{\textit{Your research question is about your problem's condition; its significance follows from your problem's cost or consequence}.}

%
What differentiates practical and conceptual problems is the nature of those conditions and costs/consequences.

The condition of a practical problem can be any state of affairs in the world that troubles you or, better, your readers: a traffic jam, foreign competition, a disease we can't effectively treat.

The cost of a practical problem is always some tangible effect we don't like: inconvenience, expense, pain, even death.

\fbox{\textit{Practical problems are often a matter of perspective}}: if my company's products are outselling yours, that's a problem for you but not for me.

%
\textit{The condition of a conceptual problem is always some version of not knowing or understanding something}.

\textit{A conceptual problem does not have a tangible cost but a consequence}.

This consequence is a particular kind of \textit{ignorance: a lack of understanding that keeps us from understanding something else that is even more significant}.

Put another way, because we haven't answered 1 question, we can't answer another that is more important.

%
\fbox{In short, practical problems concern what we should \textit{do}; conceptual problems concern what we should \textit{think}.}

Practical problems are most common in the professional world; conceptual problems are most common in academe.

\subsection{Distinguishing Pure \& Applied Research}
We call research \textit{pure} when it addresses a conceptual problem that does not have any direct practical consequences, when it only improves the understanding of a community of researchers.

We call research \textit{applied} when it addresses a conceptual problem that does have practical consequences.

You can tell whether research is pure or applied by considering the significance of your project: \textit{is it about understanding or doing?}
\begin{enumerate}
    \item \textit{Topic}: I am studying how readings from the Hubble telescope differ from readings for the same stars measured by earthbound telescopes
    \begin{itemize}
        \item[2.] \textit{Question}: because I want to find out how much the atmosphere distorts measurements of electromagnetic radiation
        \begin{itemize}
            \item[3.] \textit{Practical Significance}: so that astronomers can use data from earthbound telescopes to \textit{measure} more accurately the density of electromagnetic radiation.
        \end{itemize}
    \end{itemize}
\end{enumerate}
Applied research is common in academic fields such as business, engineering, and medicine and in companies and government agencies that do research to understand what must be known before they can solve a problem.

%
Some new researchers may be uneasy with pure research because the consequence of a conceptual problem - not knowing something - seems so abstract.

Since they are not yet part of a community that cares deeply about understanding its part of the world, they feel that their findings aren't good for much.

So they try to show the importance of their conceptual answer by cobbling onto it an implausible practical use:
\begin{enumerate}
    \item \textit{Topic}: I am studying differences among 19th-century versions of the Alamo story
    \begin{itemize}
        \item[2.] \textit{Question}: because I want to find out how politicians used stories of such events to shape public opinion
        \begin{itemize}
            \item[3.] \textit{Practical Significance}: in order to protect ourselves from unscrupulous politicians.
        \end{itemize}
    \end{itemize}
\end{enumerate}
Most readers will find the link between this research question and its asserted significance a stretch.

But for researchers in American history, the question does not need to have practical significance.

As the term \textit{pure} suggests, many researchers value the pursuit of knowledge ``for its own sake'' as a reflection of humanity's highest calling to know more.

%
So if you are doing academic research, resist the urge to turn a conceptual problem into a practical one - unless you've specifically been asked to do so.

You are unlikely to solve any genuine practical problem in a course project.

And in any case, most academic researchers see their mission not as fixing the problems of the world but as understanding them better (which may or may not lead to fixing them).

\section{Propose a Working Hypothesis}

\section{Build a Storyboard to Plan \& Guide Your Work}

\section{Join or Organize a Writing Group}

%------------------------------------------------------------------------------%

\chapter{Finding Useful Sources}

\section{3 Kinds of Sources \& Their Uses}

\section{Search for Sources Systematically}

\section{Evaluate Sources for Relevance \& Reliability}

\section{Look beyond the Usual Kinds of References}

\section{Record Your Sources Fully, Accurately, \& Appropriately}

%------------------------------------------------------------------------------%

\chapter{Engaging Your Sources}

\section{Read Generously to Understand, Then Critically to Engage}

\section{Take Notes Systematically}

\section{Take Useful Notes}

\section{Review Your Progress}

\section{Manage Moments of Normal Anxiety}

%------------------------------------------------------------------------------%

\chapter{Constructing Your Argument}

\section{What a Research Argument Is \& Is Not}

\section{Build Your Argument around Answers to Readers' Questions}

\section{Turn Your Working Hypothesis into a Claim}

\section{Assemble the Elements of Your Argument}

\section{Prefer Arguments Based on Evidence to Arguments Based on Warrants}

\section{Assemble an Argument}

%------------------------------------------------------------------------------%

\chapter{Planning a 1st Draft}

\section{Avoid Unhelpful Plans}

\section{Create a Plan That Meets Your Readers' Needs}

\section{File Away Leftovers}

%------------------------------------------------------------------------------%

\chapter{Drafting Your Paper}

\section{Draft in the Way That Feels Most Comfortable}

\section{Develop Effective Writing Habits}

\section{Keep Yourself on Track through Headings \& Key Terms}

\section{Quote, Paraphrase, \& Summarize Appropriately}

\section{Integrate Quotations in Your Text}

\section{Use Footnotes \& Endnotes Judiciously}

\section{Show How Complex or Detailed Evidence Is Relevant}

\section{Be Open to Surprises}

\section{Guard against Inadvertent Plagiarism}

\section{Guard against Inappropriate Assistance}

\section{Work Through Chronic Procrastination \& Writer's Block}

%------------------------------------------------------------------------------%

\chapter{Presenting Evidence in Tables \& Figures}

\section{Choose Verbal or Visual Representations of Your Data}

\section{Choose the Most Effective Graphic}

\section{Design Tables \& Figures}

\section{Communicate Data Ethically}

%------------------------------------------------------------------------------%

\chapter{Revising Your Draft}

\section{Check for Blind Spots in Your Argument}

\section{Check Your Introduction, Conclusion, \& Claim}

\section{Make Sure the Body of Your Report Is Coherent}

\section{Check Your Paragraphs}

\section{Let Your Draft Cool, Then Paraphrase It}

%------------------------------------------------------------------------------%

\chapter{Writing Your Final Introduction \& Conclusion}

\section{Draft Your Final Introduction}

\section{Draft Your Final Conclusion}

\section{Write Your Title Last}

%------------------------------------------------------------------------------%

\chapter{Revising Sentences}

\section{Focus on the 1st 7 or 8 Words of a Sentence}

\section{Diagnose What You Read}

\section{Choose the Right Word}

\section{Polish It Up}

\section{Give It Up \& Turn It In}

%------------------------------------------------------------------------------%

\chapter{Learning from Comments on Your Paper}

\section{2 Kinds of Feedback: Advice \& Data}

\section{Find General Principles in Specific Comments}

\section{Talk with Your Reader}

%------------------------------------------------------------------------------%

\chapter{Presenting Research in Alternative Forums}

\section{Plan Your Oral Presentation}

\section{Design Your Presentation to Be Listened To}

\section{Plan Your Poster Presentation}

\section{Plan Your Conference Proposal}

%------------------------------------------------------------------------------%

\chapter{On the Spirit of Research}

%------------------------------------------------------------------------------%

\part{Source Citation}

%------------------------------------------------------------------------------%

\chapter{General Introduction to Citation Practices}

\section{Reasons for Citing Your Sources}

\section{The Requirements of Citation}

\section{2 Citation Styles}

\section{Electronic Sources}

\section{Preparation of Citations}

\section{Citation Management Tools}

%------------------------------------------------------------------------------%

\chapter{Notes-Bibliography Style: The Basic Form}

\section{Basic Patterns}

\section{Bibliographies}

\section{Notes}

\section{Short Forms for Notes}

%------------------------------------------------------------------------------%

\chapter{Notes-Bibliography Style: Citing Specific Types of Sources}

\section{Books}

\section{Journal Articles}

\section{Magazine Articles}

\section{Newspaper Articles}

\section{Websites, Blogs, \& Social Media}

\section{Interviews \& Personal Communications}

\section{Papers, Lectures, \& Manuscript Collections}

\section{Older Works \& Sacred Works}

\section{Reference Works \& Secondary Citations}

\section{Sources in the Visual \& Performing Arts}

\section{Public Documents}

%------------------------------------------------------------------------------%

\chapter{Author-Date Style: The Basic Form}

\section{Basic Patterns}

\section{Reference Lists}

\section{Parenthetical Citations}

%------------------------------------------------------------------------------%

\chapter{Author-Date Style: Citing Specific Types of Sources}

\section{Books}

\section{Journal Articles}

\section{Magazine Articles}

\section{Newspaper Articles}

\section{Websites, Blogs, \& Social Media}

\section{Interviews \& Personal Communications}

\section{Papers, Lectures, \& Manuscript Collections}

\section{Older Works \& Sacred Works}

\section{Reference Works \& Secondary Citations}

\section{Sources in the Visual \& Performing Arts}

\section{Public Documents}

%------------------------------------------------------------------------------%

\part{Style}

%------------------------------------------------------------------------------%

\chapter{Spelling}

\section{Plurals}

\section{Possessives}

\section{Compounds \& Words Formed with Prefixes}

\section{Line Breaks}

%------------------------------------------------------------------------------%

\chapter{Punctuation}

\section{Periods}

\section{Commas}

\section{Semicolons}

\section{Colons}

\section{Question Marks}

\section{Exclamation Points}

\section{Hyphens \& Dashes}

\section{Parentheses \& Brackets}

\section{Slashes}

\section{Quotation Marks}

\section{Apostrophes}

\section{Multiple Punctuation Marks}

%------------------------------------------------------------------------------%

\chapter{Names, Special Terms, \& Titles of Works}

\section{Names}

\section{Special Terms}

\section{Titles of Works}

%------------------------------------------------------------------------------%

\chapter{Numbers}

\section{Words or Numerals?}

\section{Plurals \& Punctuation}

\section{Date Systems}

\section{Numbers Used outside the Text}

%------------------------------------------------------------------------------%

\chapter{Abbreviations}

\section{General Principles}

\section{Names \& Titles}

\section{Geographical Terms}

\section{Time \& Dates}

\section{Units of Measure}

\section{The Bible \& Other Sacred Works}

\section{Abbreviations in Citations \& Other Scholarly Contexts}

%------------------------------------------------------------------------------%

\chapter{Quotations}

\section{Quoting Accurately \& Avoiding Plagiarism}

\section{Incorporating Quotations into Your Text}

\section{Modifying Quotations}

%------------------------------------------------------------------------------%

\chapter{Tables \& Figures}

\section{General Issues}

\section{Tables}

\section{Figures}

%------------------------------------------------------------------------------%

\appendix

\chapter{Paper Format \& Submission}

\section{General Format Requirements}

\section{Format Requirements for Specific Elements}

\section{File Preparation \& Submission Requirements}

%------------------------------------------------------------------------------%

\printbibliography[heading=bibintoc]
\end{document}